\section{A Historical Trace of Capacity Utilization Measurement}
\label{sec:literature_review}

\paragraph{2.1}
This section traces the capacity utilization index not as a neutral technical indicator, but as a historical product of macroeconomic statecraft and contested political economy. The metric emerged to make idle capital legible for Keynesian demand management, transformed into a sentinel for inflation control during the crisis of Fordism, and now structures neoliberal policy through technocratic output-gap benchmarks. Throughout these historical phases, the unobserved productive ceiling has remained conceptually underidentified. We trace this trajectory through four steps: the bureaucratic standardization of survey measures, the transition to output-gap governance, the theoretical debates within the post-Keynesian tradition between Sraffian and Kaleckian economists over long-run utilization, and Shaikh’s accounting-based framework.

\subsection{From Survey Measures to an Inflation Sentinel: Capacity Utilization during the Fordist Era}
\label{subsec:fordist_survey_statecraft}

\paragraph{2.2}
Before capacity utilization could serve as a policy target, early researchers had to resolve a fundamental contradiction in its definition. The mass idleness of the Great Depression forced the state to quantify unused resources, exposing a tension between engineering and economic definitions of productive capacity \citep{Burns1954, KleinDavid1958}. Engineering capacity represents the absolute physical limit of machine operation. Economic capacity, by contrast, defines the profit-maximizing level of output given prevailing wages, input costs, machine speeds, and other factors of technological development. From its inception, measuring capacity was not a neutral descriptive exercise, but an act of statecraft designed to identify the normal conditions of operation for capitalist firms, shaping economic understanding amidst the rise of American Keynesianism.

\paragraph{2.3}
During the post-war expansion, Alvin Hansen’s secular stagnation thesis warned that chronic underutilization of capital was an inherent drag on long-run growth, forcing the state to confront the productive ceiling directly \citep{Hansen1936, Hansen1941, Barber1987}. As John Hicks and Paul Samuelson codified these concerns into the IS-LM synthesis, demand management became the core paradigm of macroeconomic policy \citep{Hicks1937, Samuelson1947, Samuelson1948}. To stabilize this post-war compromise, the state apparatus required quantitative indicators to detect productive bottlenecks and effective-demand shortfalls. Capacity utilization emerged as the primary metric to render capital idleness legible, transforming abstract stagnation concerns into concrete policy triggers.

\paragraph{2.4}
The McGraw-Hill plant-and-equipment surveys standardized this legibility by asking managers to report their operating rates relative to preferred capacity benchmarks. Emerging from the private sector, these surveys served as a central information infrastructure for corporate planning and monopoly capital business strategy, facilitating coordination and forecasting before public initiatives took root \citep{BaranSweezy1988, Munroe2007}. The surveys folded heterogeneous, firm-level engineering realities into a uniform statistical reporting stream \citep{Butler1958, Phillips1963}. This private institutionalization depoliticized the capacity concept, replacing structural debates over the efficiency of capitalist production with a standardized language of business forecasting. When the Federal Reserve Board subsequently integrated these survey responses into its official Industrial Production index, it institutionalized these subjective manager estimates as the default benchmark for macroeconomic policy \citep{Board2017}. Alternative, objective physical measures—such as the workweek of capital—never secured comparable bureaucratic adoption \citep{Foss1963, TaubmanGottschalk1971}.

\paragraph{2.5}
The stagflation crisis of the 1970s dismantled the post-war demand-stabilization framework. Faced with falling profit rates and rising inflation, policymakers repurposed the capacity index: it was no longer an indicator of effective-demand shortfalls to be corrected, but a sentinel for inflation control \citep{ShaikhManiatisPetralias1999}. High capacity utilization was redefined as a warning of upcoming price pressures, prompting monetary tightening to cool the labor market. Academic and policy debates increasingly evaluated capacity measures by their ability to predict inflation, subordinating employment targets to price-level stability \citep{KleinEtAl1973}. In these debates, Alan Greenspan argued that capacity measures must reflect systemic shortages rather than administrative averages, asserting that "if the numbers do not indicate that the economy is currently at capacity, then I suggest that the numbers are wrong, and that the correct concept of capacity must reflect the situation" \citep[p.~758]{KleinEtAl1973}. This shift did not resolve the unobserved nature of the productive ceiling, but it established the technocratic foundation for modern output-gap governance.

\subsection{From Capacity Utilization to Output-Gap Governance}
\label{subsec:output_gap_governance}

\paragraph{2.6}
During the late twentieth century, the rise of global trade and financialization reduced the policy relevance of national manufacturing surveys. Central banks and international financial institutions demanded portable, rule-based indicators to monitor economic activity across national borders \citep{Jessop2002, Jessop2003, Konings2007}. This demand institutionalized the "output gap"---the difference between actual output and potential output. This transition replaced discretionary demand management with technocratic, rule-based governance models, such as inflation targeting, Taylor policy functions, and central bank independence \citep{Taylor1993, Johnson2024}. Under the New Consensus Macroeconomics, potential output was redefined as the non-inflationary limit of economic activity, serving as an anchor to enforce wage discipline through NAIRU estimates \citep{SaadFilho2018}. Yet this policy turn did not resolve the measurement problem: potential output remained a latent variable whose shape depended entirely on modeling assumptions. Like other technocratic indicators that mask political-economic conflicts behind fragile statistical choices \citep{HerndonAshPollin2014, AshBasuDube2017}, potential output is treated by policymakers as a neutral, engineering-like benchmark that depoliticizes underlying distributional conflicts.

\paragraph{2.7}
Potential output, however, is unobservable. Hence, its estimation is contingent on specific paradigms. For example, mainstream economics typically relies on statistical filters that smooth observed output trends \citep{HodrickPrescott1997, BaxterKing1999, BeveridgeNelson1981, Hamilton2018}, aggregate production functions \citep{CBO2004PotentialGDP, Alichi2015}, or structural vector autoregressions \citep{BlanchardQuah1989}. These aggregate methodologies function as a "black box," treating capacity formation as a friction-free engineering exercise. In reality, the conversion of accumulation demand into productive capacity is institutionally and distributionally mediated, shaped by conflict at the workplace and competitive pressures in the market. The identification of potential output depends on the intensity of labor and the organization of the workplace, where class struggle over work schedules, machine speed, and supervisor discipline mediates how intensively machines run. In the market, capital strikes and shifts in the reserve army of labor regulate whether this potential capacity is realized. When low unemployment strengthens labor's bargaining power, profit squeeze dynamics prompt capital to idle capacity rather than expand production.

\paragraph{2.8}
The 2008 Global Financial Crisis exposed the structural limits of this technocratic framework. Post-crisis revisions routinely flipped the estimated sign of output gaps, and downward-revised potential-output estimates were used to justify fiscal austerity and compress public investment in the European periphery \citep{Heimberger2017, Heimberger2020}. While supranational research divisions, such as those of the IMF, documented severe hysteresis and revision volatility \citep{CerraSaxena2017, CerraFatasSaxena2023}, policy surveillance accelerated the use of high-frequency nowcasting and forward-guidance models to preserve the output-gap apparatus as adequate for real-time policymaking \citep{AastveitTrovik2014, Foster2024}. This revision-prone, closure-dependent policy framework remains entrenched, providing the institutional backdrop against which heterodox debates over capacity utilization must be understood.

\subsection{Capacity-Utilization Debates in Political Economy}
\label{subsec:pe_utilization_debates}

\paragraph{2.9}
Within the debates in political economy, the "capacity utilization controversy" defines clear dividing lines between contested approaches under a central inquiry: is the long-run normal rate of utilization endogenous to demand, or is it anchored in cost structures and classical competition? The Neo-Kaleckian tradition treats the normal rate as endogenous, arguing that firms adjust their long-run operating rates in response to demand shocks, micro-level shift choices, or historical hysteresis \citep{Lavoie1995, Nikiforos2013, Setterfield2020}. In contrast, Sraffian-Keynesian economists argue that cost minimization and classical competition anchor the normal rate of utilization in the long run, with capital investment adjusting capacity to meet changes in effective demand \citep{Ciccone1986, Kurz1986}.

\paragraph{2.10}
Both heterodox positions face conceptual and empirical challenges when identifying the productive ceiling. In the Sraffian supermultiplier framework, convergence to a normal rate requires a component of aggregate demand to remain permanently autonomous, a condition that is difficult to reconcile with the debt dynamics of stock-flow consistent models \citep{Nikiforos2018}. On the other hand, the Neo-Kaleckian claim of endogenous normal utilization depends on capacity utilization exhibiting non-stationarity (or path dependency) in the long run. This empirical premise has triggered a revival of the "capacity utilization controversy" in political economy. Critics like \citet{GahnGonzalez2020} comment on Nikiforos's data, and \citet{GahnGonzalez2022} present cross-country evidence, arguing that capacity utilization is stationary and mean-reverting, suggesting that actual utilization indeed converges to a structurally determined normal rate. However, these empirical tests remain highly contested, with their results depending entirely on the choice of data and unit-root testing specifications \citep{Nikiforos2020}.

\paragraph{2.11}
By focusing primarily on manufacturing, both the FRB and AWW series fail to capture the structural shifts of modern, service-dominated economies (Haluska, 2023). Identifying the long-run dynamics of capacity formation therefore requires an alternative empirical approach that bypasses these survey-based and smoothing constraints entirely.

\subsection{Shaikh's Structural Identification of Capacity Utilization}
\label{subsec:shaikh_structural_measurement}

\paragraph{2.12}
Rather than relying on survey measures or statistical filters that assume a symmetric business cycle and enforce mean reversion around a smooth potential trend, \citet{Shaikh2016} shifts the identification of the capacity ceiling to classical profitability accounting. This methodology addresses a core conceptual limit of conventional filters: by smoothing out fluctuations, standard statistical routines treat expansions and contractions as mirror images, thereby obscuring the fundamental asymmetry of the business cycle—where slow, grinding booms are interrupted by rapid, violent contractions. This critique of symmetric filters aligns with Latin American neo-structuralist accounts by \citet{FfrenchDavis2010, FfrenchDavis2012}, who argue that deep contractions induce a persistent "recessive gap"—a demand-driven deficit that permanently scars productive capacity through hysteresis rather than mean-reverting. Both perspectives reject smoothing filters, showing that capacity utilization measurement must capture these asymmetrical paths. Shaikh's framework anchors capacity in the dynamics of capital accumulation and profitability. Using a Generalized Perpetual Inventory Method (GPIM) to construct the capital stock, he derives the potential-output path from the accounting identities linking actual profit rates to the maximum potential rate. Under this approach, capacity utilization is identified as the residual deviation of actual output from this long-run, profit-rate-determined capacity ceiling, preserving the historical and asymmetrical traces of accumulation and scrapping.

\paragraph{2.13}
This accounting-based approach is methodologically significant because of its reproducibility and portability. By relying strictly on national accounts, alongside a theoretically grounded distinction of net capital stocks (capturing valuation dynamics) versus gross capital stocks (reflecting productive assets in operation), it provides a transparent and replicable framework for independent researchers to analyze accumulation dynamics without relying on country-specific survey designs or arbitrary statistical filters. However, treating capacity utilization as a deterministic residual assumes that output and capital share a stable, long-run cointegrating vector. If this underlying econometric relation is misspecified or unstable, the resulting capacity index will be highly sensitive to minor specification choices. This chapter critically replicates Shaikh's cointegration approach to stress-test its robustness. In doing so, it serves as a necessary stepping stone: by mapping the empirical boundaries and sensitivities of Shaikh's baseline model, we lay the groundwork to overcome its limitations in subsequent chapters.
